\section{Worked examples of the exact ledger}

This chapter works through concrete examples of the exact ledger.  The
examples are not proofs of the theorem; they illustrate the algebra
that the theorem uses.

\subsection{The orbit word}

The exact prefix word of the orbit of $7$ is
\begin{equation*}
  [2,1,2,1,1,2,1,1,1,2,2,2,2,1,1,3].
\end{equation*}
Its total weight is
\begin{equation*}
  T=25,
\end{equation*}
and the final prefix identity is
\begin{equation*}
  2^{25}\cdot34177=5^{16}\cdot7+78674588089.
\end{equation*}

\subsection{Block equations}

\begin{align*}
  [2,1]\text{ from }7:&&
  2^3\cdot23&=184=5^2\cdot7+9,\\
  [1,2]\text{ from }9:&&
  2^3\cdot29&=232=5^2\cdot9+7,\\
  [1,1,2]\text{ from }29:&&
  2^4\cdot229&=3664=5^3\cdot29+39.
\end{align*}

\subsection{PMI sum}

For $w=[2,1,2]$ from $7$,
\begin{equation*}
  \sum_{j=0}^{2}2^{-m_j}
    =1+\frac45+\frac8{25}
    =\frac{53}{25}
    =\frac{5\wordA(w)}{5^3}.
\end{equation*}
All three proper prefixes have positive margin.

\subsection{Why closure is restrictive}

Consider the word
\begin{equation*}
  w=[2,2,3].
\end{equation*}
Its molecule is
\begin{equation*}
  \wordA(w)=25+5\cdot4+2^{2+2}=25+20+16=61,
\end{equation*}
and its total weight is $7$.  If it were closed from a state $m$, then
\begin{equation*}
  2^7\,m=5^3\,m+61,
\end{equation*}
so
\begin{equation*}
  128m=125m+61,
  \qquad
  3m=61,
\end{equation*}
which has no integer solution.  The word satisfies the weight
inequality $T>3\log_2 5$ but still cannot close.

\subsection{A longer example}

For
\begin{equation*}
  w=[3,3,2,2],
\end{equation*}
the molecule is
\begin{equation*}
  \wordA(w)=125+200+320+256=901,
\end{equation*}
and the total weight is $10$.  Closedness would require
\begin{equation*}
  2^{10}\,m=5^4\,m+901,
\end{equation*}
so
\begin{equation*}
  1024m=625m+901,
  \qquad
  399m=901,
\end{equation*}
again with no integer solution.

\subsection{Insufficient weight}

For
\begin{equation*}
  w=[2,2,2,2,3],
\end{equation*}
the molecule is
\begin{equation*}
  \wordA(w)=625+500+400+320+256=2101,
\end{equation*}
and the total weight is $11$.  Since
\begin{equation*}
  11<5\log_2 5\approx11.6096,
\end{equation*}
closedness is impossible already by the weight inequality: the
left-hand side $2^{11}m$ is smaller than $5^5m$ for every positive
$m$.

\subsection{The lesson}

These examples show why the cycle equation is restrictive: the
weights must satisfy a strict margin inequality, and the molecule
must also divide the weight difference
\begin{equation*}
  \wordA(w)=\left(2^T-5^P\right)m.
\end{equation*}
The proof uses this exact equation to extract every constraint on a
cycle word.

\subsection{Status}

\begin{center}
\small
\begin{tabular}{@{}p{0.56\textwidth}p{0.36\textwidth}@{}}
\toprule
Item & Status \\
\midrule
orbit prefix identities & compiled \\
block equations & math proven \\
PMI sums & compiled \\
closure examples & exact arithmetic, not proof claims \\
\bottomrule
\end{tabular}
\end{center}
